% from ega4-errata-addenda-list3-fr.tex
\oldpage[IV]{361}
$\operatorname{div}(t_y^n)$, ce qui n'est autre que la condition
2\textdegree{} de {\hyperlink{ega4.c04.cor.21-4-10}{(21.4.10)}},
puisque $X_y$ n'a qu'un seul point maximal.
\end{proof}

\begin{remark}[21.4.12]
\label{IV.21.4.12}
Si dans {\hyperlink{ega4.c04.cor.21-4-11}{(21.4.11)}} on suppose
seulement que les $X_y$ sont réduits, il faut ajouter à la condition de
l'énoncé que les images réciproques de $D'$ dans chacun des ensembles
$\operatorname{Div}(\operatorname{Spec}(\sh{O}_{X,x}))$, où $x$
parcourt l'ensemble des points maximaux de $X_y$, sont toutes de la
forme $\operatorname{div}((t'_x)^n)$ où $t'_x$ est une uniformisante
de $\sh{O}_{X,x}$, l'entier $n$ étant le même pour tous ces points.
\end{remark}

\begin{corollary}[21.4.13]
\label{IV.21.4.13}
Supposons vérifiées les conditions de
{\hyperlink{ega4.c04.cor.21-4-11}{(21.4.11)}}. Pour tout
$\sh{O}_X$-Module inversible $\sh{L}$ tel que, pour tout point maximal
$\eta$ de $Y$, l'image réciproque
\[
 \sh{L}_\eta=\sh{L}\otimes_{\sh{O}_X}\sh{O}_{X_\eta}
\]
soit isomorphe à $\sh{O}_{X_\eta}$, il existe un
$\sh{O}_Y$-Module inversible $\sh{M}$ et un isomorphisme de
$f^*(\sh{M})$ sur $\sh{L}$.
\end{corollary}

\begin{proof}
Il suffit de montrer qu'il existe une section méromorphe régulière $s$
de $\sh{L}$ au-dessus de $X$ telle que le diviseur $D'$ correspondant
à $(\sh{L},s)$ soit tel que $D'_x=0$ en tout point $x$ d'une fibre
$X_\eta$ d'un point maximal $\eta$ de $Y$. On appliquera alors
{\hyperlink{ega4.c04.cor.21-4-11}{(21.4.11)}} à $D'$, et la conclusion
résultera de {\hyperlink{ega4.c04.def.21-4-2}{(21.4.2)}}.

Notons que les points de $\operatorname{Ass}(\sh{O}_X)$ sont au-dessus
des points maximaux de $Y$ {\hyperref[canon-index.ega4.c04.prop.3-3-1]{(3.3.1)}};
comme $Y$ est réduit et localement noethérien, il y a un ouvert dense
$V$ de $Y$, réunion d'une famille disjointe d'ouverts affines contenant
chacun un seul point maximal de $Y$, et $f^{-1}(V)$ est
schématiquement dense dans $X$
{\hyperref[canon-index.ega4.c04.prop.5-10-2]{(5.10.2)}}; on peut donc se borner au
cas où $Y$ est affine et intègre. Puisque, pour l'unique point maximal
$\eta$ de $Y$, $X_\eta$ est alors intègre,
$\operatorname{Ass}(\sh{O}_X)$ est réduit à un seul point $z$, qui est
le point générique de $X$. Soit $s_\eta$ la section de $\sh{L}_\eta$
au-dessus de $X_\eta$ qui correspond par isomorphie à la section unité
de $\sh{O}_{X_\eta}$, et soient $U$ un voisinage ouvert affine d'un
point $x\in X_\eta$ dans $X$ et $s$ une section de $\sh{L}$ au-dessus
de $U$ ayant même germe que $s_\eta$ au point $x$. Comme $U$ est
schématiquement dense dans $X$
({\hyperlink{ega4.c04.statement.20-2-13}{20.2.13}}, (iv)), $s$
s'identifie à une section méromorphe régulière bien déterminée de
$\sh{L}$ au-dessus de $X$, et il est clair
({\hyperlink{ega4.c04.prop.20-2-11}{20.2.11}} et
{\hyperlink{ega4.c04.prop.20-3-5}{20.3.5}}) que $s$ induit sur $X_\eta$
une section méromorphe régulière de $\sh{L}_\eta$ égale à $s_\eta$, et
répond donc à la question.
\end{proof}

\bigskip
\begin{center}
\small
1967. --- Imprimerie des Presses Universitaires de France. --- Vendôme
(France)\\
ÉDIT. N\textsuperscript{o} 29 470\qquad
IMPRIMÉ EN FRANCE\qquad
IMP. N\textsuperscript{o} 20 298
\end{center}
